\documentclass[11pt]{article}
\usepackage[margin=1in]{geometry}
\usepackage{amsmath,amssymb,amsthm}
\newtheorem{proposition}{Proposition}
\newcommand{\E}{\mathbb E}
\newcommand{\Prb}{\mathbb P}
\newcommand{\tr}{\operatorname{tr}}
\newcommand{\diag}{\operatorname{diag}}
\title{Isolation and mean-square approximation of the mixed target coefficient}
\author{Project research note}
\date{18 November 2026}
\begin{document}
\maketitle

Let $G=E(\mathbb F_p)$ have size $N$, with the odd prime field and
supplied group size of the project. Fix $T\in G$ with $2T\ne O$ before
drawing the seeds. In slot $i$ retain both affine points above each
abscissa satisfying $(\alpha_i x+\beta_i\bmod M)<\tau_i$, with
$M=2^b$, uniform odd $\alpha_i$, uniform $\beta_i$, independent choices
within and across slots, and dyadic $1\le\tau_i<M$. Write $f_i$ for the
point indicator, $F_i$ for its support, and
\[
 H=[T]\prod_{i=1}^m f_i,\qquad \rho_i=\tau_i/M,\qquad
 \nu_i=N\rho_i,\qquad \lambda=\frac{\prod_i\nu_i}{N}.
\]
The product is group convolution, so $H$ counts ordered tuples. This
is the coefficient that survives polarization in the second-digit note.
Let $X$ count its tuples satisfying the original generic chord guards.
We obtain constant-probability isolation at narrower widths, an
estimate computable from the individual slot sizes, and a limitation
on approximating the target-specific part on average over actual seeds.
The joint parity evaluation cost remains open.

\section{Isolation at each prescribed target}

For one hash, two different abscissas are both accepted with probability
zero when they agree modulo $\tau$, and exactly $\rho^2$ otherwise.
Indeed, condition on the first hashed value. The difference of hashed
values is uniform on the residues of the same 2-adic valuation $v$ as
the original difference. If the first value is accepted and
$2^v<\tau$, exactly $\tau/2^{v+1}$ of the $M/2^{v+1}$ possible
second values lie in the window. If $2^v\ge\tau$ there are none.
A single abscissa is accepted with probability $\rho$.

\begin{proposition}
For $m\ge2$ put $\kappa=|G[2]|\le4$. Then
\begin{equation}
 \E[H(H-1)]\ \le\ D :=
 \lambda^2\prod_i(1+2/\nu_i)
 +\lambda\,\kappa(2^m-2)/N . \label{eq:factorial}
\end{equation}
For $N\ge30m$ the probability of a unique tuple which also passes
the guards is at least $\lambda(1-3m/N)-D$.
\end{proposition}
\begin{proof}
Pair distinct tuples with sum $T$ and classify them by the set $S$
of slots in which their abscissas agree. If $S$ is proper, there are
at most $N^{m-1}$ choices for the first tuple. There are at most
$2^{|S|}$ sign choices on $S$ and $N^{m-|S|-1}$ choices elsewhere,
since the last unspecified point is forced. The joint retention
probability is at most $\prod_{i\in S}\rho_i\prod_{i\notin S}\rho_i^2$.
Summing over proper sets gives at most the first term of
(\ref{eq:factorial}).

If all abscissas agree, the second tuple is obtained by sign changes.
Use as the flip set $J$ the coordinates where the points actually
differ. It is nonempty, and flipping all entries would give $-T=T$.
The sum on $J$ lies in $G[2]$. For a fixed $J$ and a fixed such sum
$U$, the blocks $J$ and $J^c$ have prescribed sums $U$ and $T-U$.
Together they allow at most $N^{m-2}$ first tuples. The second tuple
is then fixed and retention costs only $\prod_i\rho_i$. Summing gives
the second term. Finally, the targetwise generic-fiber bound in the
source manuscript (its coverage lemma) gives
$\E X\ge(1-3m/N)\lambda$, since every generic tuple is retained with
probability $\prod_i\rho_i$. Pointwise
$1_{\{X=H=1\}}\ge X-H(H-1)$, proving the claim.
\end{proof}

Widths can be chosen publicly before drawing any hashes. Take
\[
 s=\left\lfloor\log_2\big((2/3)M^m/N^{m-1}\big)\right\rfloor
\]
and distribute $s$ as evenly as possible among the $m$ integer
exponents $\log_2\tau_i$. This is admissible in the project's
asymptotic regime, since each exponent is $b/m+O(1)$ and $m\ge2$.
It gives $1/3<\lambda\le2/3$. If $m/\min_i\nu_i=o(1)$ and
$2^m/N=o(1)$, the bound becomes
\[
     \Prb(X=H=1)\ \ge\ \lambda(1-\lambda)-o(1)
                       \ \ge\ 2/9-o(1).
\]
Both conditions hold for the growing polylogarithmic arities in the
project. Thus a single parity bit is odd on a constant mass of
instances carrying a unique guarded tuple. At $T=O$ on an odd
prime-order curve all affine tuples occur in opposite pairs, so this
extension to the identity would be false.

There is also a consequence at the original equal paid widths.
Cauchy--Schwarz and $X^2\le H^2$ give
\[
 \Prb(X>0)\ge
 \frac{(1-3m/N)^2}
 {\lambda^{-1}+\prod_i(1+2/\nu_i)+\kappa(2^m-2)/(N\lambda)} ,
\]
using $\E H\le\lambda$. At $\nu_i\ge2N^{1/m}$,
$m/\min_i\nu_i=o(1)$ and growing $m$ in this regime, this tends to
one. This pointwise average over seeds does not assert simultaneous
coverage of every target by one seed.

\paragraph{What a parity evaluator would buy.}
Suppose exact parity were available on each dyadic subwindow of the
chosen slots. Starting at an odd count, bisect one window in its
hashed coordinate and query the parity of its lower half. Exactly one
half has odd count. The upper half is another legal mask after shifting
$\beta_i$ by minus the half width. Repeat until each width is one.
Each retained abscissa then has at most two points. Check the at most
$2^m$ sign choices against the target and guards. An isolated guarded
tuple always survives. The number of parity calls including the first
is at most $1+\sum_i\log_2\tau_i$. The adaptive queries need correct
answers on the restrictions encountered. An average success guarantee
on freshly drawn full-size masks alone does not justify this reduction.
An odd count alone certifies an ordinary tuple. The guard check is
necessary before accepting a witness.

\section{How much of the mixed coefficient can individual slots explain?}

For this section assume $G$ itself has odd prime order $r=N$, and $T\ne O$.
The trivial Fourier character gives the estimate
$P=\prod_i|F_i|/r$, which costs the individual slot enumerations only.
Put $R=H-P$ and $n=(r-1)/2$. All expectations below are at a fixed
target under the original independent seed law, including both
multipliers and shifts. No distribution on curve encodings is used.

\begin{proposition}
Split the slots into any two nonempty fixed sets $L,K$. For arbitrary
real functions of their respective seeds, put
$Z=P+\sum_{a=1}^q u_a(L)v_a(K)$. Then
\begin{align}
 \E (H-P)^2 &\le \lambda(r-2)/r,\label{eq:upper}\\
 \E (H-Z)^2 &\ge
 \frac{\lambda}{r}\left[r-2-4\sum_i(1+n\rho_i)-4q\right]_+ .
 \label{eq:rank}
\end{align}
The bound with $q=0$ also holds if an arbitrary sum of functions,
each omitting at least one slot, is added to $P$.
\end{proposition}
\begin{proof}
Index rows by one character from each nontrivial opposite pair and
columns by one point $A$ from each affine opposite pair (equivalently
by its abscissa). The real $n$-by-$n$ matrix
$W_{j,A}=\chi_j(A)+\chi_j(-A)$ satisfies
$WW^\top=rI-2J$ and $W{\bf1}=-{\bf1}$ by character orthogonality.
Here $J$ is the all-ones matrix. In slot $i$ let $z_i$ be the column
of accepted-abscissa bits, and let $B_i$ be the abscissa-to-residue
incidence matrix modulo $\tau_i$. The two-abscissa law gives
\[
 \E z_i z_i^\top=\rho_i I+\rho_i^2(J-B_iB_i^\top).
\]
Define $A_i=\E (Wz_i)(Wz_i)^\top/\nu_i$. Then
\[
 I-A_i=\frac{2-\rho_i}{r}J+
             \frac{\rho_i}{r}W B_iB_i^\top W^\top .
\]
In particular $0\preceq A_i\preceq I$ and
$\tr(I-A_i)\le 1+n\rho_i$, since
$\|W\|_{\rm op}^2\le r$ and $\|B_i\|_F^2=n$. For the centered
feature $Wz_i-\E Wz_i=Wz_i+\rho_i{\bf1}$ the same proof replaces
$2-\rho_i$ by $2$ and gives the same trace bound.

For a nonempty set $S$ let $U_S$ have columns equal to the products
of the respective normalized Fourier features, indexed by the
characters, on the probability space of the seeds in $S$.
Independence gives $U_S^*U_S=A_S=\mathop{\circ}_{i\in S}A_i$
(entrywise product). Positivity of entrywise products of positive
semidefinite matrices follows by using tensor products of Gram
vectors. Applying it to $(I-A)\circ B$ bounds $A\circ B$ above by
$I\circ B\preceq I$ when both factors are positive contractions.
The diagonal product inequality $1-\prod x_i\le\sum(1-x_i)$ for
$0\le x_i\le1$ gives
$\tr(I-A_S)\le\sum_{i\in S}(1+n\rho_i)$.

Fourier inversion with opposite characters paired expresses the
residual kernel as $(\sqrt{\prod_i\nu_i}/r)U_L D_T U_K^*$, with
$D_T$ diagonal of entries $\chi_j(T)+\chi_j(-T)$.
Here $\|D_T\|_{\rm op}\le2$ and $\tr D_T^2=r-2$ since $2T\ne O$.
Squared kernel norm is mean squared error on the two probability
spaces. After removing the common squared scale $\lambda/r$, its
value is $\tr(D_T A_L D_T A_K)$. It is at most $\tr D_T^2$ and
at least $\tr D_T^2-4\tr(I-A_L)-4\tr(I-A_K)$. For the latter,
expand around the identity and observe that
$\tr(D_T(I-A_L)D_T(I-A_K))\ge0$. Every squared singular value in
this normalization is at most $4$. Projecting onto the column space
of any rank-$q$ approximation can remove at most $4q$ units of
squared norm. This proves (\ref{eq:rank}).

For the last assertion project $R$ with
$\prod_i(I-\E_i)$, where $\E_i$ averages only the seed of slot $i$.
These commuting orthogonal projections annihilate every term which
omits a slot. The remaining Fourier expression replaces all features
by their centered versions, so its squared norm has the same lower
bound. This proof allows sums over different omitted slots.
\end{proof}

At the balanced tuned or paid widths,
$m/r+\sum_i\rho_i=o(1)$. For $q=o(r)$ the mean-square error is
therefore at least $(1-o(1))\lambda$, while slot sizes alone attain
at most $\lambda$. The actual mean of $H$ is
$\prod_i\rho_i\,((r-1)^m-(-1)^m)/r$ for $T\ne O$, obtained by
convolving the full affine mask $1_G-1_{\{O\}}$. Thus at paid widths
the slot sizes give small root-mean-square error relative to the
mean, of order $\lambda^{-1/2}$, even though the absolute
root-mean-square error is of order $\sqrt{\lambda}$.

\section{Predicting the parity bit directly}
Replace $2/3$ by $1/2$ in the balanced width construction. This gives
$1/4<\lambda\le1/2$. Write $V=H\bmod 2$ and interpret a predictor
$Y$ as a matrix of output bits on a fixed nontrivial seed split.
\begin{proposition}
In the balanced asymptotic regime above, if the real rank of $Y$
is $q=o(r)$, then
\[
    \Prb(Y\ne V)\ \ge\ \Prb(V=1)-o(1)
             \ \ge\ 3/16-o(1).
\]
Always predicting even attains the first probability on the right.
The same comparison with the constant-zero predictor holds for
any Boolean event contained in $\{H>0\}$.
\end{proposition}
\begin{proof}
The proof of (\ref{eq:rank}) gives an operator norm bound on the
residual, $\|R\|_{\rm op}\le2\sqrt{\lambda/r}$ in the probability
weights of the seeds. Put $u=\E Y$. The sum of the singular values
of the weighted matrix $Y$ is at most $\sqrt q\,\|Y\|_F=\sqrt{qu}$,
by Cauchy--Schwarz. Pairing a singular-value expansion of $Y$ with
$R$ gives $|\E(RY)|\le2\sqrt{\lambda q u/r}$.

Put $p_0=\E P=\lambda(1-1/r)^m$ and $\sigma_P^2=\operatorname{Var}P$.
Each slot size has mean $\mu_i=(r-1)\rho_i$ and variance at most
$2\mu_i$, since each abscissa contributes two points and distinct
abscissa indicators have nonpositive covariance. Independence gives
\[
   \sigma_P^2\le p_0^2\left(\prod_i(1+2/\mu_i)-1\right)=o(1).
\]
Thus $\E(HY)\le p_0u+(\sigma_P+2\sqrt{\lambda q/r})\sqrt{u}$.
For any event with indicator $V'\le H$ this yields the finite bound
\[
 \Prb(Y\ne V')-\Prb(V'=1)
  \ \ge\ (1-2p_0)u
       -2(\sigma_P+2\sqrt{\lambda q/r})\sqrt{u}.
\]
Since $p_0\le1/2$, the maximal improvement tends to zero. For parity,
the isolation result with these widths gives
$\Prb(V=1)\ge\lambda(1-\lambda)-o(1)\ge3/16-o(1)$.
\end{proof}

This prediction result uses the real rank of the output bit itself.
Thresholding a low-rank real score may yield a high-rank bit matrix.
The result makes no running-time claim for algorithms that process
the slots jointly or compute over a finite field.

\section{Consequence for the time bound}
Since $\Prb(H>0)\le\E H\le\lambda$, constant coverage requires
$\prod_i\nu_i=\Omega(N)$. The isolation construction attains that
scale and keeps $\log_2(2\sum_i\tau_i)=b/m+O(\log m+1)$.
The existing conditional cost for a joint query has log scale
$mb/\sqrt{\log_2 b}$. Paying for the factor base and for these queries
in the project ledger still gives the existing conditional estimate
\[
 \log_2 \E T_{\rm solve}\ \le\
 (26.2933775156+o(1))\frac{b}{(\log_2 b)^{1/4}}
\]
under the existing finite-packet assumption A.1 and the inherited
contraction hypotheses. The constant is carried over from the earlier
short-hash construction. Isolation supplies a simpler success criterion.
The exponent would change only with a saving elsewhere in the ledger,
such as a cheaper joint parity evaluator that also handles the
restricted windows. The rank and projection bounds price specified
forms of approximation and give no general running-time lower bound.

\paragraph{Finite results and reproduction.}
Run \texttt{python -B research/mixed\_coefficient/verify.py} with NumPy.
The checks enumerate the full seed distributions by merging identical
masks and preserving their integer multiplicities. They cover the
order-19 and order-31 prime curves and an order-8 curve with four
rational two-torsion points including the identity. On the four-slot
order-19 fixture they check all seven nontrivial seed splits up to
interchanging sides. Direct bisection recovered all 42 isolated
guarded target instances in the chosen restriction fixtures.

On $y^2=x^3+44x+1$ over $\mathbb F_{65521}$, direct point counting
certifies the prime order $r=65419$. The following are evaluations of
the inequalities, with no sampling of seeds on this curve. Three
width-32 slots give $\lambda=0.49821632$ to eight decimal places,
singleton probability above $0.20027$, and a mean-square error of at
least $0.99074\lambda$ for $P$ plus any 100 separated real products.
At widths $(16,32,32)$ the Boolean parity error for real rank at most
100 exceeds $0.1135$. For this last number minimize the finite parity
bound over $\sqrt u$, bounding the maximal possible gain by
$(\sigma_P+2\sqrt{\lambda q/r})^2/(1-2p_0)$, and use the isolation
bound for $\Prb(V=1)$.

The project inputs used here are the coverage
lemma in \texttt{compressed\_relation\_search.tex}, the cost ledger in
\texttt{research/short\_hash\_count/count.tex}, and the coefficient
identity in \texttt{research/second\_digit/digit.tex}.
\end{document}
